% Generated by `taskgen course examples` from number_theory.fermat_vs_miller_rabin (seed 5).
% Do not edit: change examples.yaml and regenerate.
\begin{example}[тест Миллера--Рабина для числа Кармайкла $561$]\label{ex:miller-rabin-561}
Убедиться, что число $N = 561$ проходит тест Ферма по основанию $a = 2$. Выполнить для $N$ и $a$ тест Миллера-Рабина. Если найден нетривиальный квадратный корень из единицы по модулю $N$, найти с его помощью нетривиальный делитель $N$.

\par\noindent\textit{Решение.} $N - 1 = 560 = 2^{4} \cdot 35$. Вычисляем $x_0 = 2^{35} \bmod N$ и $x_{i} = x_{i-1}^2 \bmod N$, $i = 1, \ldots, 4$ (тогда $x_{4} = a^{N-1}$). Для $x_0$ --- бинарный метод слева направо, $d = (100011)_2$, $e$ --- обработанная часть показателя:\[\begin{array}{c|rrrrr}e & 1 & 2 & 4 & 8 & 17 \\ \hline 2^e \bmod N & 2 & 4 & 16 & 256 & 359\end{array}\]\[\begin{array}{c|r}e & 35 \\ \hline 2^e \bmod N & 263\end{array}\]Последовательность $x_i$:\[\begin{array}{c|rrrrr}i & 0 & 1 & 2 & 3 & 4 \\ \hline x_i & 263 & 166 & 67 & 1 & 1\end{array}\]$x_{4} = 2^{N-1} \equiv 1 \pmod N$: тест Ферма пройден, хотя $N = 3 \cdot 11 \cdot 17$ составное ($N$ --- число Кармайкла). Тест Миллера-Рабина: $x_0 \ne 1$ и $-1$ не встречается среди $x_0, \ldots, x_{3}$, значит, $a = 2$ --- \textbf{свидетель Миллера}, $N$ составное. При этом $x = x_{2} = 67$: $x^2 \equiv 1$, $x \not\equiv \pm 1 \pmod N$ --- нетривиальный корень из $1$. Так как $N \mid (x - 1)(x + 1)$, но $N$ не делит ни один из сомножителей, \[\begin{gathered}\gcd(x - 1, N) = \gcd(66, 561) = 33,\\ \gcd(x + 1, N) = \gcd(68, 561) = 17,\end{gathered}\] $N = 33 \cdot 17$.
\end{example}
